\section{Related Work}
\label{related_work_sec}
Parallel graph algorithms have been extensively studied due to their importance in large-scale data analysis. In this section, we review existing approaches to parallel breadth-first search (BFS), wait-free algorithm design, and connected component detection in graphs.

Several parallel BFS algorithms have been proposed for both shared-memory and distributed-memory systems. Level-synchronous BFS is a widely used approach in shared-memory systems, where all vertices at a given level are processed in parallel \cite{bader_bfs},\cite{mizell_bfs}.

Bader and Madduri \cite{bader_bfs} present a fine-grained parallelization of the level synchronous BFS algorithm designed for shared-memory systems. Their implementation emphasizes load balancing and low-overhead synchronization, achieving strong performance across a wide range of graph structures. Given that their code is publicly available, we include a direct comparison by running both their implementation and ours on the same hardware setup.

Mizell and Maschhoff \cite{mizell_bfs} later build on this work with performance improvements for a larger multithreaded system. Xia and Prasanna \cite{xia_bfs} propose synchronization reducing optimizations for multicore architectures, including an adaptive barrier to dynamically adjust thread participation, leading to improved performance. Leiserson and Schardl \cite{charles_bfs} introduce a “bag” data structure in place of a shared queue, leveraging the Cilk++ runtime for effective parallelization. While these approaches apply different optimizations and are evaluated on varied hardware and graph families, they collectively highlight the diversity of strategies for optimizing parallel BFS.

In distributed-memory settings, BFS is often implemented using message passing via MPI. For instance, Nongmeikapam \cite{aditya_bfs} presents a parallel BFS implementation using MPI with an adjacency matrix representation. While effective for small graphs, this approach is limited by the high space complexity of the adjacency matrix, making it impractical for large, sparse graphs commonly encountered in real-world applications. Scarpazza et al.'s \cite{scarpazza_bfs} BFS for the Cell/B.E. processor, though multicore, resembles distributed BFS strategies based on explicit partitioning and edge aggregation. Yoo et al. \cite{yoo_bfs} implemented a distributed BFS for BlueGene/L using 2D graph partitioning to limit communication to $\sqrt{p}$ processors, achieving scalability up to 32,000 processors. However, this assumes regular degree distributions and suffers under skewed graphs, with computation time increasing significantly under weak scaling. Cong et al. \cite{cong_bfs} explored graph algorithms in the PGAS model, and Edmonds et al. \cite{edmond_bfs} introduced the first hybrid-parallel 1D BFS using active messages.

Our work differs from the above in that we introduce a wait-free parallel BFS algorithm for shared-memory systems—an approach that, to the best of our knowledge, has not yet been explored. This design eliminates the need for global synchronization or locking mechanisms, allowing all threads to progress independently while ensuring correctness and improving scalability.

A disjoint-set union (DSU) data structure is commonly used in serial algorithms for computing connected components. The core idea extends naturally to parallel settings, and one well-known approach is the wait-free parallel union-find algorithm proposed by Anderson and Woll \cite{anderson_dsu}. Their method enables concurrent union and find operations without locking, making it suitable for high-performance parallel implementations. Randomization is also a widely used technique in this context. The random-mate (RM) algorithm, introduced by Reif and J H \cite{reif_cc}, assigns each vertex the role of either parent or child based on a fair coin toss, after which each child selects a parent from among its neighbors.

In our work, we compare against a parallel DSU-based connected components algorithm, evaluating its performance alongside our proposed methods. This allows us to assess the effectiveness of our approach in relation to a well-established parallel technique.
